\section{Exact spatial templates without a prescribed raster}
\label{sec:r15-spatial-calculus}
This new spherical profile supplies coordinate geometry and direct primitive
evaluation to the hinge engine. It is distinct from the inherited horizontal
chart $\rho=\log(H/r_0)$, $\theta=\operatorname{atan2}(N,E)$ and from the
quantized WANT.W64 address. An expression-template table stores complete
functions and parameters. It is not a table of sampled distances whose
interpolant is silently declared exact.

\subsection{Spherical logarithmic coordinates retain the Euclidean metric}
For a fixed positive reference length $r_0$, let $\rho\in\R$, polar angle
$0<\theta<\pi$, and azimuth $\varphi$ modulo $2\pi$. In fixed Cartesian axes,
\begin{equation}
 r=r_0e^\rho,\qquad
 x=o+r e_r,\qquad
 e_r=(\sin\theta\cos\varphi,\sin\theta\sin\varphi,\cos\theta).
 \label{eq:r15-log-spherical}
\end{equation}
The orthonormal tangents are
$e_\theta=(\cos\theta\cos\varphi,\cos\theta\sin\varphi,-\sin\theta)$
and $e_\varphi=(-\sin\varphi,\cos\varphi,0)$. Differentiation gives the
metric and volume element
\begin{equation}
 ds^2=r^2(d\rho^2+d\theta^2+\sin^2\theta\,d\varphi^2),
 \qquad dV=r^3\sin\theta\,d\rho\,d\theta\,d\varphi.
 \label{eq:r15-spherical-metric}
\end{equation}
Thus equal coordinate increments do not have equal physical length. The
standard orthogonal-coordinate gradient/divergence construction is given in
\href{https://web.mit.edu/8.Math/www/lectures/lec14/lec14.pdf}
{SC1: MIT Mathematical Methods, lecture 14}; substitution $r=r_0e^\rho$
gives the following explicit bodies for a smooth scalar $f$:
\begin{align}
 \nabla f&=\frac1r\left(e_r f_\rho+e_\theta f_\theta
                  +e_\varphi\frac{f_\varphi}{\sin\theta}\right),\nonumber\\
 \nabla^2f&=\frac1{r^2}\left(f_{\rho\rho}+f_\rho+
             f_{\theta\theta}+\cot\theta f_\theta+
             \csc^2\theta f_{\varphi\varphi}\right).
 \label{eq:r15-spherical-differential}
\end{align}
These factors are required if a selected gravity, diffusion or field equation
is expressed in this chart. Removing them changes the physical equation.
The origin is outside the logarithmic chart; the polar-axis and angular-seam
coordinate ambiguities use an explicit alternate chart or absent-chart status.
They are not physical holes. Full azimuth winding is retained separately from
a periodic address.

For twice differentiable coordinate functions of time and a translating
origin $o(t)$, with the Cartesian axes held fixed, direct kinematics are
\begin{align}
 \dot x-\dot o&=r(\dot\rho e_r+\dot\theta e_\theta+
                                \sin\theta\dot\varphi e_\varphi),\nonumber\\
 \ddot x-\ddot o&=r(a_r e_r+a_\theta e_\theta+a_\varphi e_\varphi),\nonumber\\
 a_r&=\ddot\rho+\dot\rho^2-\dot\theta^2-\sin^2\theta\dot\varphi^2,\nonumber\\
 a_\theta&=\ddot\theta+2\dot\rho\dot\theta
                          -\sin\theta\cos\theta\dot\varphi^2,\nonumber\\
 a_\varphi&=\sin\theta\ddot\varphi
       +2\dot\rho\sin\theta\dot\varphi+2\cos\theta\dot\theta\dot\varphi.
 \label{eq:r15-spherical-kinematics}
\end{align}
These are derivatives of the same seeded trajectory, rather than finite
differences of rendered points. Rotating axes additionally require R14's
frame-velocity and acceleration terms. A specified trajectory supplies
kinematics; its physical realization still needs the corresponding force and
energy model. Exact expression differentiation does not identify an unknown
trajectory from an uncalibrated observation.

\subsection{Distance, defining function and occupancy remain different outputs}
For a declared solid $A$ with boundary $\partial A$, its signed Euclidean
distance is the distance to $\partial A$, negative inside, positive outside
and zero on the boundary. Where differentiable away from exceptional closest-
point sets, its gradient has unit norm. In the chart above this becomes
\begin{equation}
 f_\rho^2+f_\theta^2+\csc^2\theta f_\varphi^2=r^2.
 \label{eq:r15-sdf-eikonal}
\end{equation}
The metric-weighted equation is a necessary local property; boundary data,
sign and the selected global distance definition are still required. The
unweighted expression $f_\rho^2+f_\theta^2+f_\varphi^2=1$ is not the Euclidean
SDF condition in logarithmic spherical coordinates.

A sphere has the exact direct SDF $f(x)=\|x-a\|-R$. Its squared guard in
equation~\eqref{eq:r15-sphere-guard} has the same sign but different units and
gradient. A normalized plane has $f(x)=n\cdot x-c$, $\|n\|=1$.
The retained segment and finite-cone templates likewise provide explicit
piecewise distance bodies under their own domains. An absolute-value pyramid
inequality is a valid defining predicate when its region is intended; it is
not automatically a Euclidean distance. A scalar sphere formula cannot be
replaced by an unspecified angular LUT and an assertion of identical shape.

Occupancy is $[f<0]$ with its declared boundary policy. Boolean combinations
of defining functions retain their predicate logic. For example,
$\min(f_A,f_B)<0$ describes a union, but inside overlapping solids that
minimum need not be the distance to the union's exterior boundary. Therefore
the engine exposes whether a returned scalar is a true distance, a distance
bound, a defining guard, or only a Boolean. A proof requiring a Lipschitz
constant uses the actual selected function, not its informal SDF label.

\subsection{Rigid and uniform-scale transport is exactly solvable}
Let $f_0$ be a signed distance in material coordinates. For a translation
$a(t)$, rotation $R(t)\in SO(3)$ and uniform scale $\lambda(t)>0$, set
\begin{equation}
 x=a+\lambda R X,\qquad X=R^T(x-a)/\lambda,
 \qquad f_t(x)=\lambda f_0(X).
 \label{eq:r15-similarity-sdf}
\end{equation}
For every pair $X,Y$, the transform multiplies their Euclidean separation by
$\lambda$. Taking the infimum over the original boundary proves that
\eqref{eq:r15-similarity-sdf} is exactly the transformed signed distance.
On its differentiable domain, $\nabla_x f_t=R\nabla_Xf_0$, preserving unit
normal length.

Write $\dot R=[\omega]_\times R$ in world components. The velocity of a
material point, expressed at its current position, is
\begin{equation}
 v(x,t)=\dot a+\omega\times(x-a)+(\dot\lambda/\lambda)(x-a).
 \label{eq:r15-material-velocity}
\end{equation}
Differentiating $f_t(a+\lambda RX)=\lambda f_0(X)$ at fixed $X$ gives
\begin{equation}
 \partial_t f_t+v\cdot\nabla_x f_t=(\dot\lambda/\lambda)f_t.
 \label{eq:r15-sdf-material-rate}
\end{equation}
Rigid transport has $\lambda=1$ and zero material derivative. On the moving
boundary $f_t=0$, the same relation gives its normal velocity. Corners and
medial sets keep their nondifferentiability status; a unit normal is not
invented where the selected derivative has no value.

For an arbitrary diffeomorphism $\Phi$, the pullback guard
$g_t=f_0\circ\Phi^{-1}$ preserves the transformed boundary and inside sign,
but
\begin{equation}
 \nabla_x g_t=(D\Phi)^{-T}\nabla_X f_0
 \label{eq:r15-deformed-normal}
\end{equation}
generally has nonunit length. Dividing by this local length does not in
general recover the global Euclidean distance. If the full transformation is
globally bi-Lipschitz with constants $0<\ell\le L$, the actual transformed
distance satisfies the useful bound
\begin{equation}
 \ell|f_0(X)|\le |f_t(\Phi(X))|\le L|f_0(X)|.
 \label{eq:r15-deformed-bound}
\end{equation}
This follows by taking the boundary infimum of the two distance inequalities.
Local Jacobian values alone do not establish those global constants on an
arbitrary domain. A deformation or its inverse supplied only through an ODE
retains that solution/inversion obligation.

\subsection{Direct template evaluation has a precise computational scope}
A template-table entry contains an identifier, complete expression body,
parameters, domains, frame/time binding and declared output type. At supplied
$x,t$ the selected sphere, plane, affine-coordinate or admitted algebraic
guard can be evaluated directly. Their derivatives can be supplied as exact
bodies. No spatial raster or marching search is required for that selected
evaluation. This is the useful meaning of a direct or ``query-free'' primitive
profile; it does not remove the supplied point/time or the work of evaluating
its expression.

Finding a first ray hit instead asks for a selected root of
$f(o+\tau d)=0$. Finding a general nearest feature asks for a minimum with its
domain and selection conditions. A closed primitive may supply a finite
formula; an arbitrary composite, deformation or trajectory may require root
selection, a search, or an unresolved result. Exact encoding does not turn
these different tasks into one constant-cost lookup. A finite dictionary of
templates is not a dictionary containing every possible continuous shape.

\subsection{PSI excludes a certified descendant set, not future events}
Let $N$ identify a branch and its declared descendants $\mathcal D_N$ over
the time interval $I$. Here descendant points include every point of a
represented primitive's extent, not only one origin per instance. Its bound
must assert
\begin{equation}
 \forall D\in\mathcal D_N,\ \forall t\in I:
       \|P_D(t)-P_N(t)\|\le R_N(t).
 \label{eq:r15-descendant-bound}
\end{equation}
Primitive extent, transformed children, motion and any represented uncertainty
must all be included. A center-only bound is insufficient for a finite solid.
For a 1-Lipschitz query guard $f_t$ with excluded side $f_t>0$, the condition
\begin{equation}
 f_t(P_N(t))>R_N(t)
 \label{eq:r15-psi}
\end{equation}
implies $f_t(P_D(t))>0$ for every enclosed point: subtract the distance bound
from the center value. For a whole-interval exclusion the condition must hold
throughout $I$, or follow from a common interval lower bound. Equality is not
excluded. For Lipschitz constant $L_f$, replace the right side by $L_fR_N$.
If evaluation returns $\underline f_N\le f_t(P_N)$, a certified decision
uses $\underline f_N>L_fR_N$, including all additional margins in compatible
units.

One concrete slab version is
\begin{equation}
 |n\cdot P_N-c|>w/2+R_N+d_N+\delta_N,\qquad \|n\|=1,
 \label{eq:r15-psi-slab}
\end{equation}
where $R_N$ encloses projected descendants, $d_N$ bounds their projected
deformation and $\delta_N$ bounds remaining projected uncertainty. It is the
retained corrected PSI proof: every descendant, not merely an undeformed
parent hull, is on the excluded side. Monotone-looking branch indices or a
positive depth counter cannot replace the quantified enclosure.

A finite AABB hierarchy realizes another enclosure. Each parent contains all
child bounds; leaf bounds include their actual points/primitive extents.
After motion, changed leaves are refitted and their ancestors recomputed, or
a previously certified all-time bound is used. Geometry and bounds carry the
same epoch before pruning; mixed old/new bounds do not establish enclosure.
A self-referential return that changes the root or transforms descendants
starts a new bound evaluation unless the old certificate explicitly covered
that return. An infinite descendant family requires a proved finite bound or
cannot be pruned by that finite-radius certificate.

Pruning here removes only a proved irrelevant geometric evaluation for a
specified query and epoch. It does not skip accepted hinge events, JK updates,
clock progression, winding or future seam returns. Diffraction, thermal
coupling, force or sensor contributions require their R14 error/energy
obligations before they are omitted from a physical model. Keeping those
obligations preserves the distinction between a fast spatial engine and a
proved physical simulation.
